\documentclass[11pt,a4paper]{article}

\usepackage[utf8]{inputenc}
\usepackage[T1]{fontenc}
\usepackage{lmodern}
\usepackage[italian]{babel}
\usepackage{amsmath,amssymb,mathtools}
\usepackage[expansion=false]{microtype}
\usepackage[a4paper,margin=2.6cm]{geometry}
\usepackage{setspace}
\usepackage[most]{tcolorbox}
\usepackage{fontspec}
\usepackage[italian, bidi=basic, provide=*]{babel}
\babelprovide[import, onchar=ids fonts]{italian}
\babelfont{rm}{Noto Sans}
\usepackage{tikz}
\usepackage{amsmath}
\usetikzlibrary{arrows.meta, calc, backgrounds}

\setstretch{1.12}
\setlength{\parindent}{1.2em}
\setlength{\parskip}{0.3em}

% Notazione
\newcommand{\dd}{\mathrm{d}}
\newcommand{\epz}{\varepsilon_0}
\newcommand{\Ev}{\vec{E}}
\newcommand{\Fv}{\vec{F}}
\newcommand{\uv}{\hat{u}}
\newcommand{\kc}{\dfrac{1}{4\pi\epz}}

% Riquadro per i richiami (titolo senza grassetto)
\newtcolorbox{richiamo}[1]{
  enhanced, breakable,
  colback=gray!5, colframe=gray!55!black,
  boxrule=0.6pt, arc=2pt,
  left=8pt, right=8pt, top=6pt, bottom=6pt,
  fonttitle=\normalfont\large, coltitle=white,
  title={#1}
}

\begin{document}

\begin{center}
{\LARGE Il campo elettrostatico nel vuoto\\[0.4ex]
dalla legge di Coulomb alle distribuzioni continue di carica}\\[1.2ex]
{\large Dispensa di Fisica 2 per la preparazione all'esame}
\end{center}

\vspace{1em}

\noindent
Il percorso prende avvio dalla legge di Coulomb, che descrive la forza tra due cariche puntiformi ferme come azione diretta a distanza, e approda alla nozione di campo elettrostatico, che sposta l'attenzione dalla coppia di corpi alla regione di spazio che circonda la sorgente. Fissata la definizione operativa, il principio di sovrapposizione permette di costruire il campo di un insieme qualunque di cariche: somme vettoriali per i sistemi discreti, integrali di linea, di superficie e di volume per le distribuzioni continue. Seguono i casi canonici, cioè il dipolo, il filo indefinito, l'anello, il disco, il piano indefinito, il doppio strato e lo strato sferico, scelti perché in ciascuno la simmetria riduce il calcolo vettoriale a un solo integrale scalare. Una regolarità attraversa tutti gli esempi: a distanza grande rispetto alle dimensioni della sorgente, il campo coincide con quello di una carica puntiforme, mentre a distanza piccola emerge la geometria della distribuzione. La dispensa è scritta in modo da poter essere letta senza altri testi; i richiami di Meccanica e di Analisi sono raccolti negli appositi riquadri.

\medskip

\begin{richiamo}{Richiamo Matematico/Fisica 1: forze, equilibrio, forza elastica ed energia}
Una forza è una grandezza vettoriale, misurata in newton, e la sua azione su un punto materiale di massa $m$ è regolata dalla relazione $\Fv_{\mathrm{tot}}=m\vec a$, dove $\Fv_{\mathrm{tot}}$ è la somma vettoriale di tutte le forze applicate e $\vec a$ l'accelerazione. Il terzo principio della dinamica afferma che, se un corpo A esercita su un corpo B la forza $\Fv_{AB}$, il corpo B esercita su A la forza $\Fv_{BA}=-\Fv_{AB}$, con la stessa retta d'azione. Se $\Fv_{\mathrm{tot}}=\vec 0$, l'accelerazione è nulla: il punto resta fermo oppure procede di moto rettilineo uniforme, e in entrambi i casi la velocità è costante. Se invece $\Fv_{\mathrm{tot}}$ è costante e non nulla, l'accelerazione è costante e il moto, lungo una retta, è uniformemente accelerato: $x(t)=x_0+v_0t+\tfrac12 a t^2$.

La forza elastica $F = -kx$ ($k$ in $\mathrm{N/m}$) è sempre diretta verso la posizione di equilibrio $x = 0$ e conduce all'equazione $\ddot{x} = -(k/m)x$, le cui soluzioni sono oscillazioni armoniche di pulsazione $\omega = \sqrt{k/m}$ (in $\mathrm{rad/s}$) e periodo $T = 2\pi/\omega$. Ogni volta che una forza risulta proporzionale e opposta allo spostamento, il moto è un oscillatore armonico, qualunque sia l'origine fisica della forza.

In contesti unidimensionali, le forze che dipendono solo dalla posizione risultano conservative, garantendo la conservazione dell'energia meccanica $E_m = \frac{1}{2}mv^2 + U$ ($U$ energia potenziale, in joule); l'energia potenziale si converte in cinetica e viceversa, e il punto di inversione del moto si trova dove $v = 0$. In presenza di attrito parte dell'energia meccanica si dissipa; in tutti gli esempi che seguono l'attrito è trascurato per dichiarazione esplicita.

%La forza elastica $F=-kx$ ($k$ in $\ensuremath{\mathrm{N/m}}$) è sempre diretta verso la posizione di equilibrio $x=0$ e conduce all'equazione $\ddot x=-(k/m)\,x$, le cui soluzioni sono oscillazioni armoniche di pulsazione $\omega=\sqrt{k/m}$ (in $\ensuremath{\mathrm{rad/s}}$) e periodo $T=2\pi/\omega$. Ogni volta che una forza risulta proporzionale e opposta allo spostamento, il moto è un oscillatore armonico, qualunque sia l'origine fisica della forza. Per forze che dipendono solo dalla posizione, l'energia meccanica $E_m=\tfrac12 mv^2+U$ si conserva ($U$ energia potenziale, in joule); l'energia potenziale si converte in cinetica e viceversa, e il punto di inversione del moto si trova dove $v=0$. In presenza di attrito parte dell'energia meccanica si dissipa; in tutti gli esempi che seguono l'attrito è trascurato per dichiarazione esplicita.
\end{richiamo}

\begin{richiamo}{Richiamo Matematico/Fisica 1: vettori, versori, approssimazioni e geometria}
Un vettore $\vec r$ ha modulo $r=|\vec r|$ e direzione individuata dal versore $\uv=\vec r/r$, vettore di modulo unitario e privo di dimensioni. La somma di vettori segue la regola del parallelogramma, o equivalentemente quella punta-coda; operativamente, si scompone ogni vettore lungo gli assi scelti e si sommano le componenti omonime. Se tutti i vettori della somma, una volta riportati in catena punta-coda, formano un poligono chiuso, la risultante è nulla.

Per un triangolo qualsiasi di lati $a$, $b$, $c$ vale il teorema di Carnot: $c^2=a^2+b^2-2ab\cos\gamma$, dove $\gamma$ è l'angolo opposto al lato $c$; nel caso $\gamma=\pi/2$ si ricade nel teorema di Pitagora. Per piccole quantità $\epsilon$ vale lo sviluppo di Taylor al primo ordine $(1+\epsilon)^{\alpha}\simeq 1+\alpha\epsilon$, con $\alpha$ reale: tale approssimazione è il principale strumento per ottenere il comportamento di un campo a grande distanza. Le aree utili sono: circonferenza di raggio $s$, di lunghezza $2\pi s$; corona circolare di raggio interno $s$ e spessore infinitesimo $\dd s$, di area $\dd S=2\pi s\,\dd s$; superficie sferica di raggio $R$, di area $4\pi R^2$; fascia sferica compresa tra le colatitudini $\theta$ e $\theta+\dd\theta$, di area $\dd S=2\pi R^2\sin\theta\,\dd\theta$.

Un integrale definito $\int_a^b f(x)\,\dd x$ è il limite di una somma di contributi infinitesimi $f(x)\,\dd x$ e si calcola tramite una primitiva $F$ come $F(b)-F(a)$. Quando una grandezza vettoriale è sommata su una distribuzione continua, l'integrale si scrive per componenti; la simmetria della configurazione consente spesso di dimostrare in anticipo che alcune componenti si annullano a coppie. Gli angoli sono sempre espressi in radianti: un radiante è l'angolo che sottende un arco di lunghezza pari al raggio, per cui $\pi/2$ corrisponde a $90^{\circ}$ e $2\pi$ a $360^{\circ}$; in Fisica 2 i gradi non compaiono nelle formule.
\end{richiamo}

La forza elettrica tra corpi carichi fermi è la prima interazione di cui si conosca una forma quantitativa esatta, e il suo accertamento è sperimentale. Lo strumento è la bilancia di torsione, concettualmente affine al pendolo con cui Cavendish misurò l'attrazione gravitazionale tra masse: due sfere cariche, abbastanza piccole e distanti da poter essere schematizzate come punti, producono una torsione del filo proporzionale alla forza, e la ripetizione della misura a distanze e cariche differenti mostra che il modulo della forza è proporzionale al prodotto delle cariche e inversamente proporzionale al quadrato della distanza. La legge è dunque nata per cariche puntiformi; per corpi estesi si procede suddividendoli in elementi abbastanza piccoli da essere trattati come puntiformi, ed è questo il ruolo delle distribuzioni continue trattate più avanti.

Siano $q$ e $q_0$ due cariche puntiformi, la prima posta in un punto $P$ e la seconda in un punto $P'$. Il vettore $\vec r$ va da $P$ a $P'$, ha modulo $r$ e individua il versore $\uv_r=\vec r/r$. La forza che $q$ esercita su $q_0$ vale
\[
\Fv_{q\to q_0}=\frac{1}{4\pi\epz}\,\frac{q\,q_0}{r^2}\,\uv_r ,
\qquad
\epz=\ensuremath{8,8542\cdot 10^{-12}\,\mathrm{C^2/(N\,m^2)}}.
\]

\begin{figure}[htbp]
  \centering
  \begin{tikzpicture}[>=Stealth, font=\small]

    % Stili personalizzati per elementi grafici
    \tikzset{
      gridline/.style={draw=gray!25, line width=0.3pt, dashed},
      badge/.style={draw=gray!40, fill=gray!10, rounded corners=3pt, inner xsep=8pt, inner ysep=4pt, font=\small\sffamily},
      forcearrow/.style={line width=1.3pt, ->},
      posarrow/.style={line width=0.8pt, ->, color=black!80}
    }

    % ==========================================
    % CASO 1: Repulsivo (q q' > 0)
    % ==========================================
    \begin{scope}[shift={(0,0)}]
      % Griglia tecnica di sfondo
      \draw[gridline] (-2.0,-1.1) grid (7.8,1.1);

      % Anelli concentrici (rappresentazione infografica di campo/potenziale)
      \foreach \r in {0.25, 0.45, 0.65} {
        \draw[gray!45, line width=0.4pt, dashed] (0,0) circle (\r);
        \draw[gray!45, line width=0.4pt, dashed] (4.2,0) circle (\r);
      }

      % Vettore Posizione r
      \draw[posarrow] (0,0) -- (4.2,0) node[midway, below=3pt] {$\vec{r}$};

      % Vettori Forza F' e F con prolungamento tratteggiato
      \draw[forcearrow, color=orange!90!red] (0,0) -- (-1.4,0) node[above=2pt] {$\vec{F}'$};
      \draw[dashed, color=orange!90!red, line width=0.8pt] (-1.4,0) -- (-1.8,0);

      \draw[forcearrow, color=orange!90!red] (4.2,0) -- (5.6,0) node[above=2pt] {$\vec{F}$};
      \draw[dashed, color=orange!90!red, line width=0.8pt] (5.6,0) -- (6.0,0);

      % Particelle / Cariche
      \filldraw[fill=orange!80!red, draw=black!80, line width=0.6pt] (0,0) circle (3.2pt);
      \filldraw[fill=purple!70!blue, draw=black!80, line width=0.6pt] (4.2,0) circle (3.2pt);

      % Etichette P, q, P', q'
      \node[above=4pt, font=\small\bfseries] at (0,0) {$P$};
      \node[below=4pt] at (0,0) {$q$};
      \node[above=4pt, font=\small\bfseries] at (4.2,0) {$P'$};
      \node[below=4pt] at (4.2,0) {$q'$};

      % Badge identificativo della condizione
      \node[badge] at (6.8,0) {$q\,q' > 0$};
    \end{scope}

    % ==========================================
    % CASO 2: Attrattivo (q q' < 0)
    % ==========================================
    \begin{scope}[shift={(0,-2.5)}]
      % Griglia tecnica di sfondo
      \draw[gridline] (-2.0,-1.1) grid (7.8,1.1);

      % Anelli concentrici
      \foreach \r in {0.25, 0.45, 0.65} {
        \draw[teal!40, line width=0.4pt, densely dotted] (0,0) circle (\r);
        \draw[teal!40, line width=0.4pt, densely dotted] (4.2,0) circle (\r);
      }

      % Vettore Posizione r
      \draw[posarrow] (0,0) -- (4.2,0) node[midway, below=3pt] {$\vec{r}$};

      % Vettori Forza F' e F (elevati per evitare sovrapposizioni)
      \draw[forcearrow, color=teal!80!black] (0,0.3) -- (1.4,0.3) node[midway, above=1pt] {$\vec{F}'$};
      \draw[dotted, color=teal!60, line width=0.6pt] (0,0) -- (0,0.3);

      \draw[forcearrow, color=teal!80!black] (4.2,0.3) -- (2.8,0.3) node[midway, above=1pt] {$\vec{F}$};
      \draw[dotted, color=teal!60, line width=0.6pt] (4.2,0) -- (4.2,0.3);

      % Particelle / Cariche
      \filldraw[fill=cyan!80!black, draw=black!80, line width=0.6pt] (0,0) circle (3.2pt);
      \filldraw[fill=blue!80!black, draw=black!80, line width=0.6pt] (4.2,0) circle (3.2pt);

      % Etichette P, q, P', q'
      \node[above=4pt, font=\small\bfseries] at (0,0) {$P$};
      \node[below=4pt] at (0,0) {$q$};
      \node[above=4pt, font=\small\bfseries] at (4.2,0) {$P'$};
      \node[below=4pt] at (4.2,0) {$q'$};

      % Badge identificativo della condizione
      \node[badge] at (6.8,0) {$q\,q' < 0$};
    \end{scope}

  \end{tikzpicture}
  \caption{Interazione elettrostatica tra due cariche puntiformi $q$ e $q'$ nei casi repulsivo ($q\,q' > 0$) ed attrattivo ($q\,q' < 0$).}
  \label{fig:legge_coulomb_infografica}
\end{figure}

Nel modulo compaiono $q$ e $q_0$, i valori algebrici delle cariche in coulomb, e $r$, la distanza tra i punti in metri; $\epz$ è la costante dielettrica del vuoto, grandezza che caratterizza il mezzo e non le cariche; il fattore $1/(4\pi\epz)$, spesso indicato con $k$, vale circa $\ensuremath{9\cdot 10^{9}\,\mathrm{N\,m^2/C^2}}$, e la verifica dimensionale è immediata, perché $\ensuremath{\mathrm{N\,m^2/C^2}}\cdot\ensuremath{\mathrm{C^2/m^2}}=\ensuremath{\mathrm{N}}$. Il versore $\uv_r$ punta da $q$ verso $q_0$: se le cariche hanno lo stesso segno, il prodotto $qq_0$ è positivo e la forza su $q_0$ è concorde con $\uv_r$, quindi repulsiva; se hanno segno opposto, il prodotto è negativo e la forza è attrattiva. La forza che $q_0$ esercita su $q$ è uguale e opposta, come richiesto dal terzo principio, e giace sulla stessa retta. In aria la costante dielettrica differisce da $\epz$ in misura trascurabile, per cui la formula del vuoto vale anche per le applicazioni ordinarie.

L'unità di carica è il coulomb, simbolo $\ensuremath{\mathrm{C}}$. Va ricordato che l'unità fondamentale del Sistema Internazionale è l'ampere, mentre il coulomb è un'unità derivata: un coulomb è la carica trasportata in un secondo da una corrente di un ampere, cioè $\ensuremath{1\,\mathrm{C}}=\ensuremath{1\,\mathrm{A\,s}}$. Il carattere derivato del coulomb è spesso rovesciato per errore, dato che l'ampere viene percepito come coulomb al secondo. Il valore assoluto della carica dell'elettrone, e dunque quella del protone, è $e=\ensuremath{1,602\cdot 10^{-19}\,\mathrm{C}}$, e conviene ricordarlo, perché compare in quasi ogni stima di ordine di grandezza. Per fissare le scale, si ricorda che una presa domestica italiana eroga nominalmente $\ensuremath{230\,\mathrm{V}}$ a $\ensuremath{50\,\mathrm{Hz}}$ (nelle lezioni è citato il valore storico di $\ensuremath{220\,\mathrm{V}}$) e che gli interruttori più comuni sopportano $\ensuremath{16\,\mathrm{A}}$.

La forza di Coulomb è enorme rispetto alle scale della vita quotidiana. Due cariche di $\ensuremath{1\,\mathrm{C}}$ poste a $\ensuremath{1\,\mathrm{m}}$ di distanza si respingono con circa $\ensuremath{9\cdot 10^{9}\,\mathrm{N}}$, cioè con un peso equivalente a quello di quasi un milione di tonnellate: una situazione che i corpi macroscopici non sperimentano, perché sono globalmente neutri. Un'applicazione quotidiana si trova nel comportamento del sale da cucina in acqua. Il cloruro di sodio è un solido ionico, con ioni $\mathrm{Na}^+$ e $\mathrm{Cl}^-$ disposti in un reticolo cubico a facce centrate e separati da circa $\ensuremath{0,28\,\mathrm{nm}}$; la forza di attrazione tra una coppia di ioni, nel vuoto, vale
\[
F=\frac{1}{4\pi\epz}\frac{e^2}{r^2}\simeq \ensuremath{8,99\cdot 10^{9}}\cdot\frac{(\ensuremath{1,602\cdot 10^{-19}})^2}{(\ensuremath{0,28\cdot 10^{-9}})^2}\ \ensuremath{\mathrm{N}}\simeq \ensuremath{2,9\cdot 10^{-9}\,\mathrm{N}}.
\]
Tale valore, minuscolo in assoluto, è enorme alla scala del singolo ione, e il legame è per questo tenace. Immerso in acqua, tuttavia, il cristallo si scioglie. La costante dielettrica dell'acqua è circa ottanta volte quella del vuoto, per cui nella formula di Coulomb la quantità $\epz$ va sostituita da $\epz\epsilon_r$ con $\epsilon_r\simeq 80$, e la forza tra gli ioni scende a circa $\ensuremath{3,7\cdot 10^{-11}\,\mathrm{N}}$. Il fenomeno ha anche una ragione geometrica: la molecola $\mathrm{H_2O}$ non è lineare, i due atomi di idrogeno formano un angolo con l'ossigeno, e il baricentro della carica positiva non coincide con quello della negativa; la molecola è polare, cioè si comporta da dipolo elettrico, concetto che sarà precisato più avanti. I dipoli dell'acqua si orientano attorno agli ioni e schermano il legame, che si indebolisce fino a rompersi. Si distingue così il legame ionico dal legame covalente, nel quale gli elettroni sono messi in comune: nel diamante, dove ogni atomo di carbonio occupa il centro di un tetraedro, o nel quarzo e nei vetri silicatici, dove un atomo di silicio è legato a quattro di ossigeno, la rottura è molto più difficile, ed è questa la causa della durezza del diamante e delle proprietà isolanti, elettriche e termiche, del vetro.

Il confronto con la gravità, condotto sulla coppia di protoni, mostra quanto la forza elettrica prevalga a livello microscopico. In un nucleo due protoni distano dell'ordine di un femtometro, $\ensuremath{1\,\mathrm{fm}}=\ensuremath{10^{-15}\,\mathrm{m}}$; i prefissi nano, pico e femto valgono rispettivamente $10^{-9}$, $10^{-12}$ e $10^{-15}$. La repulsione elettrica vale
\[
F_e=\frac{1}{4\pi\epz}\frac{e^2}{r^2}\simeq \ensuremath{8,99\cdot 10^{9}}\cdot\frac{(\ensuremath{1,602\cdot 10^{-19}})^2}{(\ensuremath{10^{-15}})^2}\ \ensuremath{\mathrm{N}}\simeq \ensuremath{2,3\cdot 10^{2}\,\mathrm{N}},
\]
con un peso equivalente di circa $\ensuremath{24\,\mathrm{kg}}$ (cioè quanto servirebbe per tenere sospesa una massa di tale entità). L'attrazione gravitazionale, calcolata con la costante di gravitazione universale $G=\ensuremath{6,674\cdot 10^{-11}\,\mathrm{N\,m^2/kg^2}}$ e con la massa del protone $m_p=\ensuremath{1,673\cdot 10^{-27}\,\mathrm{kg}}$, è
\[
F_g=G\,\frac{m_p^2}{r^2}\simeq \ensuremath{6,674\cdot 10^{-11}}\cdot\frac{(\ensuremath{1,673\cdot 10^{-27}})^2}{(\ensuremath{10^{-15}})^2}\ \ensuremath{\mathrm{N}}\simeq \ensuremath{1,9\cdot 10^{-34}\,\mathrm{N}},
\]
e il rapporto $F_e/F_g$ è dell'ordine di $10^{36}$. Nel nucleo, dunque, la gravità non conta nulla, e la sola interazione elettrica condurrebbe alla dispersione immediata dei protoni. Poiché la materia esiste, si deve concludere che nel nucleo agiscono altre interazioni, la forte e la debole; i neutroni, privi di carica e quindi indifferenti alla repulsione di Coulomb, contribuiscono al legame come una sorta di collante. L'equilibrio è però delicato: al crescere del numero atomico servono sempre più neutroni, e gli elementi più pesanti sono instabili e decadono. L'uranio, attraverso una lunga catena che passa per il radon (gas radioattivo, per cui l'analisi del radon è un passo obbligato prima di costruire in locali interrati), termina nel piombo, che è stabile. L'elettrone non contribuisce in modo sostanziale a questo equilibrio, perché è lontano dal nucleo per una distanza dell'ordine di $10^4$--$10^5$ volte le dimensioni nucleari.

Con la stessa procedura si apprezza il peso della gravità per corpi ordinari. Due persone di massa $\ensuremath{70\,\mathrm{kg}}$ poste a $\ensuremath{1\,\mathrm{m}}$ si attraggono con $F=G\,m^2/r^2=\ensuremath{6,674\cdot 10^{-11}}\cdot 4900\ \ensuremath{\mathrm{N}}\simeq\ensuremath{3,3\cdot 10^{-7}\,\mathrm{N}}$, equivalenti al peso di una massa di circa $\ensuremath{33\,\mathrm{\mu g}}$. La stessa interazione diventa rilevante solo con masse planetarie: il peso di una persona di $\ensuremath{70\,\mathrm{kg}}$ sulla Terra è di circa $\ensuremath{7\cdot 10^{2}\,\mathrm{N}}$, attrazione dovuta al pianeta. L'abitudine a eseguire simili stime numeriche, con verifica dell'ordine di grandezza, rientra tra le competenze richieste in sede d'esame.

Per quanto riguarda l'impostazione dei problemi, conviene che ogni domanda sulla forza di Coulomb sia accompagnata da un disegno: si collocano la carica $q$ in $P$ e la carica $q_0$ in $P'$, si traccia la distanza $r$, si sceglie il versore $\uv_r$ lungo di essa e si rappresenta $\Fv_{q\to q_0}$ applicata in $P'$ nel verso appropriato; per semplicità si ragiona con cariche positive, e il segno si ripristina alla fine. Il principio di sovrapposizione stabilisce che, dato un insieme di cariche $q_i$ fisse, la forza risultante su $q_0$ è la somma vettoriale delle forze esercitate da ciascuna, come se le altre non ci fossero:
\[
\Fv=\sum_i \Fv_i=\sum_i \frac{1}{4\pi\epz}\frac{q_i\,q_0}{r_i^2}\,\uv_i=q_0\,\sum_i \frac{1}{4\pi\epz}\frac{q_i}{r_i^2}\,\uv_i ,
\]
dove $r_i$ e $\uv_i$ sono distanza e versore riferiti alla carica $q_i$. Il fattore $q_0$ si porta fuori dalla somma, e questa osservazione è l'origine del concetto di campo. Come esempio, si dispongano quattro cariche uguali sui vertici di un quadrato e si ponga una quinta carica nel centro: ciascuna delle prime quattro dista dal centro di mezza diagonale, per cui le quattro forze hanno uguale modulo; a due a due, quelle dei vertici opposti sono uguali e contrarie, e la risultante è nulla. La carica centrale, soggetta a $\Fv_{\mathrm{tot}}=\vec 0$, resta ferma (oppure si muove di moto rettilineo uniforme, se ha velocità iniziale non nulla).

Il campo elettrostatico nasce dall'osservazione che, nell'espressione sopra, la parte che dipende dalle sorgenti è separabile dalla carica che subisce la forza. Il campo $\Ev$ generato in un punto da un sistema di cariche ferme è definito come la forza risultante su una carica di prova positiva $q_0$ posta in quel punto, divisa per $q_0$; poiché la carica di prova perturba a sua volta la distribuzione, la definizione corretta prevede un limite:
\[
\Ev=\lim_{q_0\to 0}\frac{\Fv}{q_0}.
\]
Il campo è un vettore, misurato in $\ensuremath{\mathrm{N/C}}$ (equivalente a $\ensuremath{\mathrm{V/m}}$), e ha il significato fisico di forza per unità di carica positiva posta nel punto. Il limite $q_0\to 0$ ha una motivazione fisica: la carica di prova genera un proprio campo, e se questo non è trascurabile rispetto a quello primario, la misura modifica il fenomeno che intende rilevare. La definizione si riferisce a una qualunque regione di spazio in cui un corpo carico in quiete subisce una forza, ed è quindi indipendente dalle sorgenti, e vale anche per campi non prodotti da cariche ferme. Per questo si distingue il campo elettrostatico, generato da cariche ferme, dal campo elettrico in senso più ampio. Il campo ha esistenza oggettiva: sebbene sia misurabile soltanto tramite una sonda, la sua presenza non dipende da questa.

La ragione per cui si preferisce il campo alla forza è fisica oltre che formale. Nella legge di Coulomb, lo spostamento di $q$ modifica istantaneamente $r$ e dunque la forza su $q_0$; l'avverbio ``istantaneamente'' è però inaccettabile, perché nessun segnale può propagarsi con velocità superiore a quella della luce nel vuoto, $c\simeq\ensuremath{3\cdot 10^{8}\,\mathrm{m/s}}$. La descrizione corretta è a due stadi: la sorgente crea il campo nello spazio circostante, e il campo agisce sul corpo carico con un ritardo pari al tempo di propagazione; la legge di Coulomb, con azione diretta e istantanea, ne è il limite statico. Un vantaggio pratico è che, se si conoscono le sorgenti, il campo si determina senza alcun riferimento a una carica sonda; viceversa, se è noto il campo in ogni punto di una regione, si prevede il moto di una carica collocata in un punto qualsiasi di quella regione anche senza conoscere le sorgenti. Una domanda tipica all'orale consiste proprio nel descrivere il campo di una carica puntiforme con l'ausilio di un disegno.

Per una sorgente puntiforme $q$ in $P$ e un punto $P'$ a distanza $r$, la definizione, applicata alla legge di Coulomb, dà
\[
\Ev(P')=\frac{1}{4\pi\epz}\,\frac{q}{r^2}\,\uv_r ,
\]
dove $q$ è la carica sorgente (in coulomb), $r$ la distanza tra sorgente e punto (in metri) e $\uv_r$ il versore che va dalla sorgente al punto. Non compare alcun segno meno: per $q>0$ il campo è uscente dalla sorgente, concorde con $\uv_r$, per $q<0$ il campo è entrante, discorde. Il campo ha il punto di applicazione nel punto $P'$ considerato, direzione radiale e modulo $E=|q|/(4\pi\epz r^2)$; tutti i punti equidistanti da $q$, cioè quelli di una superficie sferica centrata sulla sorgente, hanno campo di uguale modulo. Poiché $E\propto r^{-2}$, se la distanza raddoppia il campo diventa quattro volte più piccolo, e se si dimezza diventa quattro volte più grande. Notando che la formula si distingue da quella della forza per la presenza di una sola carica: l'altra carica è, per così dire, sostituita dalla carica unitaria di prova.

Il principio di sovrapposizione si trasferisce direttamente al campo: per un sistema di $N$ cariche puntiformi, in un punto $P'$ si ha
\[
\Ev=\sum_{i=1}^{N}\Ev_i=\sum_{i=1}^{N}\frac{1}{4\pi\epz}\,\frac{q_i}{r_i^2}\,\uv_i ,
\]
dove $r_i$ è la distanza di $P'$ dalla carica $q_i$ e $\uv_i$ il versore diretto da $q_i$ verso $P'$; la somma è vettoriale, e l'errore più comune consiste nel sommare i moduli. Nei quattro vertici di un quadrato, con quattro cariche uguali, il campo nel centro è nullo: i quattro contributi hanno lo stesso modulo e a coppie sono opposti. Una carica abbandonata lì, a velocità nulla, rimane ferma, perché la forza $q_0\Ev$ è nulla. Un secondo esempio è il campo al centro di un triangolo equilatero con tre cariche uguali nei vertici. Il centro, equidistante dai vertici, è anche baricentro e ortocentro; i tre campi hanno lo stesso modulo $E$ e formano tra loro angoli di $120^{\circ}=2\pi/3$. Si fissa l'asse $y$ lungo uno dei tre vettori: gli altri due, simmetrici rispetto a tale asse, hanno componenti lungo $x$ opposte e uguali, che si cancellano, mentre le componenti lungo $y$ valgono ciascuna $E\cos(2\pi/3)=-E/2$. Il campo totale ha quindi componente $y$ pari a $E-\tfrac E2-\tfrac E2=0$, e il campo risultante è nullo. La stessa conclusione si ottiene dalla regola punta-coda: tre vettori di egual modulo e a $120^{\circ}$ formano un triangolo chiuso.

Quando la sorgente è una distribuzione continua, la somma diventa un integrale. Si suddivide la distribuzione in elementi di carica $\dd q$, ciascuno trattato come puntiforme, e si somma vettorialmente:
\[
\Ev=\frac{1}{4\pi\epz}\int \frac{\dd q}{r^2}\,\uv_r ,
\]
con $\dd q=\lambda\,\dd l$ per una distribuzione lineare di densità $\lambda$ ($\ensuremath{\mathrm{C/m}}$), $\dd q=\sigma\,\dd S$ per una superficiale di densità $\sigma$ ($\ensuremath{\mathrm{C/m^2}}$) e $\dd q=\rho\,\dd\tau$ per una volumica di densità $\rho$ ($\ensuremath{\mathrm{C/m^3}}$). In ogni integrale, $r$ e $\uv_r$ sono distanza e versore dall'elemento sorgente al punto, e variano con l'elemento. Il metodo ricorrente è il seguente: si cerca la simmetria che annulla alcune componenti, si esprime la componente superstite in funzione di una sola variabile (un angolo oppure una distanza), si integra. Il campo è sempre un vettore, e conviene scrivere il modulo solo dopo aver fissato la direzione.

\begin{richiamo}{Richiamo Matematico: sviluppi in serie per distanze grandi e angoli piccoli}
Nel confronto con una sorgente puntiforme si ricorre di frequente a due approssimazioni. La prima riguarda gli angoli: per $\theta$ in radianti e $|\theta|\ll 1$ si ha $\sin\theta\simeq\theta$ e $\cos\theta\simeq 1-\theta^2/2$. L'accuratezza è maggiore di quanto ci si aspetti: per $\theta=\pi/6\simeq 0{,}524$ si ha $\sin\theta=0{,}5$ esatto contro $0{,}524$ approssimato, con scarto di circa il $5\%$. La seconda è lo sviluppo binomiale $(1+\epsilon)^{\alpha}\simeq 1+\alpha\epsilon$ per $|\epsilon|\ll 1$: ad esempio, $(1\pm\epsilon)^{-2}\simeq 1\mp2\epsilon$ e $(1+\epsilon)^{-1/2}\simeq 1-\epsilon/2$. In entrambi i casi l'errore commesso è del secondo ordine nel parametro piccolo.
\end{richiamo}

Il dipolo elettrico è la coppia di cariche $+q$ e $-q$ poste a distanza $a$. Si definisce il momento di dipolo $\vec p=q\,\vec a$, con $\vec a$ diretto dalla carica negativa alla positiva ($\ensuremath{\mathrm{C\,m}}$). Si pone l'asse $z$ lungo $\vec a$, con $+q$ in $z=+a/2$ e $-q$ in $z=-a/2$, e si calcola il campo su due rette particolari, a distanza grande rispetto ad $a$. Su un punto dell'asse, a quota $z>a/2$, i due campi sono entrambi lungo $z$ e hanno versi opposti: la carica $+q$ è più vicina e dà il contributo maggiore, per cui
\[
E_z=\frac{q}{4\pi\epz}\left[\frac{1}{(z-a/2)^2}-\frac{1}{(z+a/2)^2}\right]
=\frac{q}{4\pi\epz z^2}\left[\Bigl(1-\frac{a}{2z}\Bigr)^{-2}-\Bigl(1+\frac{a}{2z}\Bigr)^{-2}\right].
\]
Per $z\gg a$ lo sviluppo $(1\mp\epsilon)^{-2}\simeq 1\pm2\epsilon$ con $\epsilon=a/(2z)$ dà $\bigl[(1+a/z)-(1-a/z)\bigr]=2a/z$, e quindi
\[
E_z\simeq\frac{q}{4\pi\epz z^2}\cdot\frac{2a}{z}=\frac{2p}{4\pi\epz\,z^3}.
\]
Su un punto del piano mediano, a distanza $y\gg a$ dal centro, le due cariche distano entrambe $r=\sqrt{y^2+a^2/4}$. Le componenti lungo $y$ si annullano per simmetria, mentre quelle lungo $z$ sono concordi e rivolte da $+q$ verso $-q$, ossia nel verso negativo di $z$; ciascuna vale $\dfrac{q}{4\pi\epz r^2}\cdot\dfrac{a/2}{r}$. Se ne deduce
\[
E_z=-\frac{q\,a}{4\pi\epz\,(y^2+a^2/4)^{3/2}}\simeq-\frac{p}{4\pi\epz\,y^3},
\]
dove l'ultima approssimazione trascura $a^2/4$ rispetto a $y^2$. In entrambe le direzioni, dunque, il campo del dipolo decresce come $r^{-3}$, più rapidamente di quello della carica puntiforme, perché nel complesso la carica del sistema è nulla e i contributi tendono a compensarsi; sull'asse il campo è doppio che sul piano mediano e di verso opposto.

La distribuzione rettilinea uniforme indefinita è un filo di lunghezza infinita con densità lineare costante $\lambda$. Si calcola il campo in un punto $P$ a distanza $d$ dal filo. Un elemento $\dd l$ del filo, a distanza $r$ da $P$, produce un campo $\dd E=\dfrac{\lambda\,\dd l}{4\pi\epz r^2}$ che forma un angolo $\theta$ con la perpendicolare condotta da $P$ al filo. Gli elementi simmetrici rispetto al piede della perpendicolare hanno componenti parallele al filo uguali e opposte, per cui sopravvive la sola componente radiale $\dd E_\perp=\dd E\cos\theta$. Si sceglie l'angolo $\theta\in(-\pi/2,\pi/2)$ come variabile: con la coordinata $l=d\tan\theta$ lungo il filo, misurata dal piede della perpendicolare, si ha $\dd l=d\,\dd\theta/\cos^2\theta$ e $r=d/\cos\theta$, da cui
\[
E_\perp=\int_{-\pi/2}^{\pi/2}\frac{\lambda}{4\pi\epz}\,\frac{\cos^2\theta}{d^2}\,\cos\theta\,\frac{d\,\dd\theta}{\cos^2\theta}
=\frac{\lambda}{4\pi\epz d}\int_{-\pi/2}^{\pi/2}\cos\theta\,\dd\theta
=\frac{\lambda}{2\pi\epz\,d}.
\]
Il campo è radiale, diretto verso l'esterno del filo se $\lambda>0$, e il suo modulo decresce come $1/d$ (non come $1/d^2$, poiché la sorgente non è puntiforme); le unità sono quelle consuete, perché $\lambda/\epz$ ha dimensioni di $\ensuremath{\mathrm{N/C}}\cdot\ensuremath{\mathrm{m}}$.

Il campo sull'asse dell'anello carico è l'esempio che prepara il disco. L'anello ha raggio $R$ e densità lineare $\lambda$ uniforme, per carica totale $Q=2\pi R\lambda$; il punto $P$ si trova sull'asse dell'anello, a distanza $x$ dal suo centro. Ogni elemento $\dd l$ dell'anello dista da $P$ la stessa quantità $r$, poiché $r^2=x^2+R^2$ (triangolo rettangolo con cateti $x$ e $R$). Elementi diametralmente opposti danno contributi $\dd E$ uguali in modulo, le cui componenti perpendicolari all'asse si annullano a coppie, e rimane la componente lungo l'asse. Sia $\theta$ l'angolo tra $\vec r$ e l'asse: tutti gli elementi lo vedono uguale. Si ha
\[
\dd E_x=\dd E\cos\theta=\frac{1}{4\pi\epz}\,\frac{\lambda\,\dd l}{r^2}\cos\theta ,\qquad \cos\theta=\frac{x}{r}=\frac{x}{(x^2+R^2)^{1/2}},
\]
e con la sostituzione $r^2=x^2+R^2$ si trova
\[
\dd E_x=\frac{1}{4\pi\epz}\,\frac{\lambda\,x\,\dd l}{(x^2+R^2)^{3/2}} .
\]
L'unica variabile d'integrazione è $\dd l$: $\lambda$, $x$ e $R$ sono costanti, perché la posizione di $P$ è fissata e l'anello è uniforme. Si integra su tutto l'anello, con $\oint \dd l=2\pi R$:
\[
E_x=\frac{1}{4\pi\epz}\,\frac{\lambda\,x}{(x^2+R^2)^{3/2}}\,2\pi R=\frac{1}{4\pi\epz}\,\frac{Q\,x}{(x^2+R^2)^{3/2}} ,
\]
dove si è identificata $\lambda\,2\pi R=Q$, carica totale in coulomb. Il campo è diretto lungo l'asse, uscente dal piano dell'anello se $Q>0$. Per $x=0$ si ottiene $E_x=0$, coerente con l'intuizione che al centro i contributi si compensino a coppie. Per $x\gg R$ si può trascurare $R^2$ rispetto a $x^2$, per cui $(x^2)^{3/2}=x^3$ e quindi $E_x\simeq Q/(4\pi\epz x^2)$: da lontano l'anello appare come una carica puntiforme. Vale infine un'osservazione utile per l'oscillatore: per $x\ll R$ si ha $(x^2+R^2)^{3/2}\simeq R^3$ e dunque $E_x\simeq Qx/(4\pi\epz R^3)$, un campo proporzionale a $x$. Una carica $q<0$ di massa $m$ posta vicino al centro di un anello positivo subisce la forza $F_x=qE_x=-\dfrac{|q|Q}{4\pi\epz R^3}\,x$, di tipo elastico, e compie oscillazioni armoniche di pulsazione $\omega=\sqrt{|q|Q/(4\pi\epz mR^3)}$. Poiché non è stato richiesto di memorizzare il risultato, ma di saperlo ricavare, in sede d'esame vale più il ragionamento che la formula.

Si passa ora al disco di raggio $R$, con densità superficiale uniforme $\sigma$ e carica totale $q=\sigma\pi R^2$. Esiste un procedimento diretto, basato sull'angolo solido ($\dd\Omega=\dd S\cos\alpha/r^2$, con $\alpha$ semiapertura con cui il punto vede il disco e non l'angolo generico $\theta$, come corretto a lezione), ma qui si preferisce la scomposizione in anelli, che sfrutta il risultato precedente. Si considera la corona circolare di raggio $s$ e spessore $\dd s$, di area $\dd S=2\pi s\,\dd s$ e carica $\dd q=\sigma\,2\pi s\,\dd s$; per $\dd s$ infinitesimo, essa è assimilabile a un anello. In un punto dell'asse a quota $z$, il campo infinitesimo, diretto lungo l'asse, vale
\[
\dd E_z=\frac{1}{4\pi\epz}\,\frac{z\,\dd q}{(z^2+s^2)^{3/2}}=\frac{\sigma\,z}{2\epz}\,\frac{s\,\dd s}{(z^2+s^2)^{3/2}} ,
\]
e la sola variabile è $s$, che va da $0$ a $R$. Si adotta la primitiva $\displaystyle\int\frac{s\,\dd s}{(z^2+s^2)^{3/2}}=-\frac{1}{\sqrt{z^2+s^2}}$, verificabile per derivazione, e si trova
\[
E_z=\frac{\sigma\,z}{2\epz}\left[\frac{1}{|z|}-\frac{1}{\sqrt{z^2+R^2}}\right]
=\frac{\sigma}{2\epz}\,\frac{z}{|z|}\left[1-\frac{|z|}{\sqrt{z^2+R^2}}\right].
\]
In tale espressione $\sigma$ è in $\ensuremath{\mathrm{C/m^2}}$, $z$ e $R$ in metri e $E_z$ in $\ensuremath{\mathrm{N/C}}$; il fattore $z/|z|$ vale $+1$ da una parte del disco e $-1$ dall'altra e segnala che il campo si allontana dal disco da ambo i lati (per $\sigma>0$). Per $z\gg R$ si sviluppa $|z|/\sqrt{z^2+R^2}=(1+R^2/z^2)^{-1/2}\simeq 1-R^2/(2z^2)$, per cui la parentesi tende a $R^2/(2z^2)$ e si ha $E_z\simeq\dfrac{\sigma R^2}{4\epz z^2}=\dfrac{\sigma\pi R^2}{4\pi\epz z^2}=\dfrac{q}{4\pi\epz z^2}$: da lontano il disco è una carica puntiforme. Per $z\to 0^{+}$ si ha $E_z\to\sigma/(2\epz)$; per $z\to 0^{-}$, $E_z\to-\sigma/(2\epz)$. Il campo ha quindi una discontinuità finita $\sigma/\epz$ nell'attraversamento del disco, e da ciascun lato tende a zero come $1/z^2$ per $|z|\to\infty$. Il grafico è dunque costituito da due rami simmetrici (dispari in $z$), decrescenti in modulo a partire dal valore $\sigma/(2\epz)$ in prossimità del disco: dalla formula si deve saper vedere l'andamento della grandezza fisica.

Se il raggio $R$ diventa molto grande rispetto a $z$, il termine $|z|/\sqrt{z^2+R^2}$ tende a zero e si ottiene $E=\sigma/(2\epz)$: è il campo di un piano indefinito carico con densità uniforme $\sigma$. Tale campo è uniforme in modulo (non dipende dalla distanza dal piano) e perpendicolare al piano, uscente per $\sigma>0$ da entrambi i lati: $E=+\sigma/(2\epz)$ per $z>0$ e $E=-\sigma/(2\epz)$ per $z<0$. L'indipendenza dalla distanza ha un'interpretazione geometrica: un piano infinito, guardato da lontano, resta un piano infinito, e nulla nell'aspetto della sorgente cambia al variare della distanza, mentre una sfera, un anello o un disco da lontano diventano punti. Di conseguenza, una carica di prova abbandonata vicino a un piano uniformemente carico subisce una forza costante e si muove di moto uniformemente accelerato, lungo $z$ nel verso in cui è respinta. Qualora la sorgente sia un piano positivo e la carica di prova sia negativa, la carica è attratta: se il disco ha un foro, la particella passa dall'altra parte, dove la forza ha verso opposto, e per conservazione dell'energia si arresta in un punto simmetrico a quello di partenza rispetto al disco. Il moto è periodico ma non armonico: vicino al disco, la forza ha modulo costante $|q|\sigma/(2\epz)$, con segno che inverte al passaggio, e il moto è formato da tratti uniformemente accelerati e decelerati. La somiglianza con l'oscillatore armonico, evocata a lezione, è corretta qualitativamente (la forza richiama verso la posizione centrale, e vi è conversione tra energia potenziale e cinetica), ma la proporzionalità tra forza e spostamento è realizzata propriamente nel caso dell'anello, o di una corona circolare, come mostra l'esercizio finale.

Il doppio strato è formato da due piani indefiniti paralleli, uniformemente carichi con densità $+\sigma$ e $-\sigma$. Fra i due piani i campi dei due strati sono concordi, diretti dal piano positivo verso il negativo, e si sommano: $E=\dfrac{\sigma}{2\epz}+\dfrac{\sigma}{2\epz}=\dfrac{\sigma}{\epz}$. All'esterno sono discordi e di uguale modulo, per cui $E=0$. Il risultato, utile da ricordare, descrive il campo interno di un condensatore piano ideale; gli schermi capacitivi dei telefoni, ad esempio, sfruttano proprio la variazione di tale campo quando un dito si accosta alla piastra.

Lo strato sferico uniformemente carico è una superficie sferica di raggio $R$ con densità superficiale $\sigma$ costante e carica totale $q=4\pi R^2\sigma$; non va confuso con la sfera piena, che ha comportamento diverso. Da lontano il campo si riduce a quello di una carica puntiforme. Per il calcolo esatto in un punto $P$ a distanza $x$ dal centro, sull'asse $z$, si suddivide la superficie in anelli coassiali, ciascuno individuato dalla colatitudine $\theta$ e dallo spessore $\dd\theta$. La carica dell'anello è $\dd q=\sigma\,2\pi R^2\sin\theta\,\dd\theta$ e tutti i suoi punti distano da $P$ la stessa quantità $r$. Il teorema di Carnot applicato al triangolo con vertici nel centro, in $P$ e in un punto dell'anello fornisce due relazioni: $r^2=x^2+R^2-2xR\cos\theta$, e, indicato con $\varphi$ l'angolo tra $\vec r$ e l'asse in $P$, $R^2=x^2+r^2-2xr\cos\varphi$, da cui
\[
\cos\varphi=\frac{x^2+r^2-R^2}{2xr}.
\]
La derivata della prima relazione rispetto a $\theta$ dà $2r\,\dd r=2xR\sin\theta\,\dd\theta$, cioè $\sin\theta\,\dd\theta=\dfrac{r\,\dd r}{xR}$, e la variabile può essere cambiata da $\theta$ a $r$. La componente assiale del campo dell'anello è $\dd E=\dfrac{1}{4\pi\epz}\dfrac{\dd q}{r^2}\cos\varphi$, e per sostituzione si trova
\[
\dd E=\frac{1}{4\pi\epz}\,\frac{\sigma\,2\pi R^2}{r^2}\cdot\frac{r\,\dd r}{xR}\cdot\frac{x^2+r^2-R^2}{2xr}
=\frac{\sigma R}{4\epz x^2}\left[1+\frac{x^2-R^2}{r^2}\right]\dd r .
\]
L'integrazione è su $r$, tra il valore minimo $r_{\min}$ (anello più vicino a $P$, per $\theta=0$ o $\theta=\pi$) e il massimo $r_{\max}$ (anello opposto), cioè tra $|x-R|$ e $x+R$. Poiché $\int\bigl[1+(x^2-R^2)/r^2\bigr]\dd r=r-(x^2-R^2)/r$, si distinguono due casi. Per un punto esterno ($x>R$) gli estremi sono $x-R$ e $x+R$, e il valore dell'integrale è
\[
\bigl[(x+R)-(x-R)\bigr]-(x^2-R^2)\Bigl[\frac{1}{x+R}-\frac{1}{x-R}\Bigr]=2R+2R=4R,
\]
per cui $E=\dfrac{\sigma R}{4\epz x^2}\cdot4R=\dfrac{\sigma R^2}{\epz x^2}=\dfrac{q}{4\pi\epz x^2}$. Per un punto interno ($x<R$) gli estremi sono $R-x$ e $R+x$, e si ha
\[
\bigl[(R+x)-(R-x)\bigr]-(x^2-R^2)\Bigl[\frac{1}{R+x}-\frac{1}{R-x}\Bigr]=2x-2x=0,
\]
da cui $E=0$. In sintesi, fuori dallo strato il campo è quello di una carica puntiforme $q$ collocata al centro; dentro è nullo; sulla superficie il campo ha modulo $q/(4\pi\epz R^2)$ in corrispondenza del limite esterno e $0$ in corrispondenza del limite interno, con una discontinuità $\sigma/\epz$, la stessa incontrata nell'attraversamento del disco e del piano. Ne segue che il campo assume il suo valore massimo sulla superficie. Lo stesso risultato si ricava con minor fatica tramite il teorema di Gauss, che verrà trattato in seguito.

Le linee di forza sono curve il cui tangente, in ogni punto, è parallela e concorde al campo in quel punto. Si addensano dove il campo è più intenso, non si incrociano mai (altrimenti in un punto vi sarebbero due direzioni di campo) e hanno origine nelle cariche positive e termine in quelle negative; per una carica positiva isolata sono semirette radiali uscenti, per una negativa rientranti, per un piano indefinito sono rette parallele ed equispaziate, a testimonianza dell'uniformità del campo.

Per quanto riguarda gli esami precedenti, le domande orali più frequenti su questo argomento nascono da imprecisioni concettuali. Viene spesso richiesto di definire il campo elettrico e di descrivere il campo di una carica puntiforme: si attende che sia tracciato il disegno completo, con la sorgente, il punto, la distanza, il versore e il vettore campo applicato nel punto e non lungo l'intero segmento, e che sia chiarito il motivo per cui nella definizione compaiono una carica di prova positiva e il limite per $q_0\to0$. Un altro quesito ricorrente riguarda la ragione per cui si introduce il campo accanto alla forza, e si risponde con l'argomento della propagazione a velocità finita e con la possibilità di prevedere l'effetto sul corpo carico senza conoscere le sorgenti. Sono frequenti anche le domande sui limiti: che cosa avviene al campo di un anello, di un disco o di uno strato sferico a grande distanza, che cosa accade al centro dell'anello, che cosa si può dire del campo all'interno di uno strato sferico e perché il campo del disco sia discontinuo. Le equazioni di Maxwell vanno invece conosciute a memoria, mentre per le formule di anello, disco e filo è richiesto il ragionamento che le produce, non la loro riproduzione mnemonica.

Per quanto concerne gli esercizi scritti, il tema si presenta tipicamente in tre forme. La prima è il calcolo del campo di un piccolo numero di cariche puntiformi collocate su vertici di figure regolari (quadrati, triangoli equilateri), con richiesta del campo o della forza su una carica in un punto particolare, tipicamente il centro; la risoluzione si basa su simmetria e componenti. La seconda è il calcolo del campo, o della forza su una carica di prova, lungo l'asse di un anello, di un disco o di una corona circolare, spesso seguito dal limite a grande distanza e dalla stima del periodo di oscillazione di una carica abbandonata. La terza richiede di confrontare forza elettrica e gravitazionale o di stimare ordini di grandezza per atomi, nuclei e ioni, e premia la padronanza di potenze di dieci e unità.

Gli errori più penalizzati sono di pochi tipi. Il primo riguarda il segno: il versore va sempre da sorgente a punto, e dimenticarlo inverte la direzione del campo per cariche negative. Il secondo è la somma dei moduli al posto della somma vettoriale, che manca di cogliere le cancellazioni per simmetria. Il terzo è l'uso della formula del disco con $z>0$ anche per $z<0$, con omissione del fattore $z/|z|$ e dunque del cambio di segno. Il quarto è la confusione tra $\sigma/(2\epz)$, valido per un solo piano, e $\sigma/\epz$, valido tra le piastre di un doppio strato. Il quinto è l'impiego del campo di carica puntiforme dentro uno strato sferico, dove vale zero. Vi sono poi errori tecnici: trattare $x$ come variabile d'integrazione per l'anello, mentre è costante e l'unica variabile è $\dd l$; confondere $r$ (distanza dall'elemento al punto) con il raggio $R$ dell'anello; sbagliare le unità, specie nelle conversioni da centimetri o microcoulomb; e scrivere gli angoli in gradi dentro formule che li richiedono in radianti.

L'esercizio che segue condensa questi strumenti. Una corona circolare piana, di raggio interno $a=\ensuremath{2,0\,\mathrm{cm}}$ e raggio esterno $b=\ensuremath{4,0\,\mathrm{cm}}$, è carica con densità superficiale uniforme $\sigma=\ensuremath{1,0\,\mathrm{\mu C/m^2}}$. Si chiede di determinare il campo sull'asse della corona, a quota $z$, e il periodo delle piccole oscillazioni di una particella di carica $q=-\ensuremath{1,0\,\mathrm{nC}}$ e massa $m=\ensuremath{1,0\,\mathrm{mg}}$ abbandonata sull'asse a quota $z_0\ll a$, con gravità e attrito trascurati.

La corona è la differenza di due dischi concentrici, con la stessa densità $\sigma$; per sovrapposizione, il suo campo è la differenza tra quello del disco di raggio $b$ e quello del disco di raggio $a$. In alternativa, si integra direttamente la formula del disco tra $s=a$ ed $s=b$. In entrambi i casi, per $z>0$,
\[
E_z=\frac{\sigma\,z}{2\epz}\left[\frac{1}{\sqrt{z^2+a^2}}-\frac{1}{\sqrt{z^2+b^2}}\right],
\]
espressione in cui i termini $1/|z|$ del disco si elidono, e che risulta dispari in $z$ (sul lato $z<0$ vale la stessa formula, perché il fattore $z$ cambia segno). Per $|z|\ll a$ si sviluppa al primo ordine: $1/\sqrt{z^2+a^2}\simeq\dfrac1a\bigl(1-\dfrac{z^2}{2a^2}\bigr)$, e analogamente per $b$, e i termini in $z^2$ sono trascurabili rispetto ai principali. Il campo diviene
\[
E_z\simeq\frac{\sigma}{2\epz}\Bigl(\frac1a-\frac1b\Bigr)z ,
\]
lineare nella quota, a differenza di quanto avviene per il disco pieno. Sulla particella di carica $q<0$ agisce $F_z=qE_z=-k\,z$ con
\[
k=\frac{|q|\,\sigma}{2\epz}\Bigl(\frac1a-\frac1b\Bigr)=\frac{\ensuremath{1\cdot 10^{-9}}\cdot\ensuremath{1\cdot 10^{-6}}}{2\cdot\ensuremath{8,854\cdot 10^{-12}}}\,(50-25)\ \ensuremath{\mathrm{1/m}}\simeq\ensuremath{1,4\cdot 10^{-3}\,\mathrm{N/m}},
\]
dove si è usato $1/a=\ensuremath{50\,\mathrm{1/m}}$ e $1/b=\ensuremath{25\,\mathrm{1/m}}$. La forza ha verso opposto allo spostamento perché la carica è negativa e la corona positiva, e questo è il passaggio in cui un errore di segno trasformerebbe l'oscillazione in una fuga. L'equazione del moto è $m\ddot z=-kz$, e la pulsazione e il periodo valgono
\[
\omega=\sqrt{\frac{k}{m}}=\sqrt{\frac{\ensuremath{1,41\cdot 10^{-3}}}{\ensuremath{1\cdot 10^{-6}}}}\ \ensuremath{\mathrm{rad/s}}\simeq\ensuremath{38\,\mathrm{rad/s}},\qquad
T=\frac{2\pi}{\omega}\simeq\ensuremath{0,17\,\mathrm{s}}.
\]
Occorre infine verificare la validità dell'ipotesi: lo sviluppo trascura termini relativi dell'ordine di $z_0^2/(2a^2)$, che per $z_0=\ensuremath{2\,\mathrm{mm}}$ valgono circa lo $0{,}5\%$; ne segue che, per ampiezze di qualche millimetro, il moto è armonico con ottima approssimazione, mentre per $z_0$ confrontabile con $a$ il periodo dipende dall'ampiezza e il moto non è più armonico. Il periodo trovato è indipendente dall'ampiezza (isocronismo), proprietà caratteristica dell'oscillatore armonico, e dipende dalla carica e dalla massa della particella come $\sqrt{m/|q|}$, per cui una particella dieci volte più carica oscilla con un periodo ridotto di un fattore $\sqrt{10}$. Come verifica finale, per $b\to\infty$ il termine $1/b$ scompare e si ricade nel caso del disco con foro, mentre per $b\to a$ la corona si assottiglia fino a un anello e $k\to0$ al primo ordine, coerentemente con il fatto che la densità lineare equivalente $\sigma(b-a)$ tende a zero.

\end{document}
